\section{Levantamento e especificação dos requisitos}\label{sec:requisitos}

Os requisitos foram definidos a partir dos objetivos do sistema, da revisão bibliográfica e documental, da entrevista com uma psicopedagoga (Capítulo~\ref{capitulo4}) e da análise do processo de avaliação psicopedagógica, e foram divididos em funcionais e não funcionais, conforme a classificação apresentada na Seção~\ref{sec:fund_requisitos}.

\subsection{Requisitos funcionais}

Os requisitos funcionais foram definidos considerando as principais etapas do processo de avaliação psicopedagógica contempladas pelo sistema. Entre elas estão o gerenciamento dos aprendentes, o registro de informações, o gerenciamento dos instrumentos avaliativos, a aplicação das avaliações, o acompanhamento dos registros e a elaboração de relatórios.

O sistema deve permitir a autenticação dos usuários e o controle de acesso às funcionalidades de acordo com suas permissões. Após a autenticação, o profissional deve ser capaz de cadastrar e gerenciar os dados dos aprendentes, mantendo as informações organizadas e disponíveis para consulta durante os diferentes atendimentos. 

A visibilidade dos aprendentes cadastrados deve ser restrita ao psicopedagogo responsável por cada um deles, de modo que um profissional não tenha acesso aos registros de aprendentes vinculados a outros usuários. O mesmo princípio se aplica aos instrumentos avaliativos: cada psicopedagogo visualiza somente os instrumentos que criou. O perfil administrador, responsável pela gestão das contas de usuário, não acessa os dados clínicos dos aprendentes, pois esses registros são, pelo Código de Ética do Psicopedagogo, de acesso restrito ao profissional responsável \cite{abpp2019_codigo_etica}; sua exceção limita-se aos instrumentos avaliativos, que não contêm dados de crianças. Para que um aprendente não fique sem profissional com acesso a ele quando seu responsável deixa o serviço, o administrador deve poder transferir a responsabilidade para outro psicopedagogo, sem consultar o conteúdo dos registros.

Também deve ser possível registrar informações relacionadas ao processo de avaliação, incluindo dados provenientes de entrevistas, anamnese, observações e aplicação de instrumentos. Dessa forma, os registros de um aprendente podem ser mantidos de maneira estruturada, permitindo o acompanhamento das informações ao longo do processo avaliativo.

Outro aspecto relevante é a rastreabilidade. Os registros do sistema devem conter informações suficientes para identificar quem os criou, quando foram criados e quando foram modificados. Os campos de controle presentes nas entidades (\textit{createdAt}, \textit{updatedAt} e, em parte delas, \textit{deletedAt}) e as referências ao usuário responsável por cada registro, como \textit{responsavelId} em aprendentes, \textit{autorId} em avaliações, \textit{aplicadorId} em aplicações e \textit{autorId} em relatórios, indicam quem criou cada dado e quando ele foi alterado pela última vez. Como esses campos não são suficientes para dados de saúde de crianças (Seção~\ref{sec:seguranca}), o sistema deve manter também um registro de auditoria das operações, incluindo as consultas aos dados clínicos e as operações realizadas pelo administrador.

Durante a aplicação de um instrumento, o profissional deve poder vinculá-lo a um determinado aprendente e registrar as informações obtidas no atendimento. As adaptações realizadas em um instrumento são registradas por meio de versionamento automático: quando as questões de um instrumento já aplicado são alteradas, o sistema arquiva a versão anterior e cria uma nova, mantendo cada aplicação associada à versão utilizada no momento do atendimento.

O sistema deve permitir ainda a consulta ao histórico de avaliações e registros associados ao aprendente, possibilitando que o profissional acompanhe as informações produzidas em diferentes momentos do processo avaliativo. A organização longitudinal desses registros busca facilitar a consulta e a comparação das informações produzidas ao longo do acompanhamento.

Por fim, deve ser possível elaborar devolutivas e relatórios a partir dos registros realizados no sistema, permitindo sua organização em um documento e sua posterior geração em formato PDF. O sistema atua, portanto, como instrumento de apoio à organização das informações, sem substituir a análise e a interpretação realizadas pelo profissional.

Com base nessas definições, os principais requisitos funcionais contemplados no desenvolvimento são apresentados no Quadro~\ref{qua:requisitos_funcionais}.

\begin{quadro}[H]
\centering
\small
\caption{Requisitos funcionais}
\begin{tabular}{p{2cm} p{13cm}}
\hline
\textbf{Código} & \textbf{Requisito funcional} \\
\hline
RF01 & Permitir o cadastro e a autenticação de usuários. \\
RF02 & Permitir o controle de acesso às funcionalidades do sistema de acordo com o perfil do usuário. \\
RF03 & Permitir o cadastro, edição e consulta dos dados dos aprendentes. \\
RF04 & Restringir a visualização dos aprendentes e de seus registros clínicos ao psicopedagogo responsável por cada um deles, sem exceção para o perfil administrador, a quem cabe apenas transferir a responsabilidade entre profissionais. \\
RF05 & Permitir o registro de informações relacionadas ao processo de avaliação. \\
RF06 & Permitir o cadastro e gerenciamento dos instrumentos avaliativos. \\
RF07 & Restringir a visualização dos instrumentos avaliativos ao psicopedagogo que os criou, com exceção do perfil administrador. \\
RF08 & Permitir o gerenciamento de diferentes versões dos instrumentos. \\
RF09 & Registrar adaptações realizadas nos instrumentos por meio de versionamento automático, preservando a versão anterior e associando-a às aplicações já realizadas. \\
RF10 & Permitir vincular um instrumento a um aprendente durante um atendimento. \\
RF11 & Permitir registrar os resultados e observações obtidos durante a aplicação dos instrumentos. \\
RF12 & Permitir consultar o histórico de avaliações e registros de um aprendente. \\
RF13 & Permitir elaborar e armazenar devolutivas e relatórios. \\
RF14 & Permitir gerar relatórios em formato PDF. \\
RF15 & Manter nos registros as informações de autoria e temporalidade necessárias para identificar quem criou cada dado e quando foi criado ou modificado. \\
\hline
\end{tabular}
\fonte{Elaborado pela autora (2026).}
\label{qua:requisitos_funcionais}
\end{quadro}

Os requisitos funcionais foram representados também por meio de um diagrama de casos de uso da UML \cite{omg2017_uml}, apresentado na Figura~\ref{fig:casos_de_uso}, que indica as funcionalidades disponíveis a cada perfil de usuário. Os dois atores compartilham apenas a autenticação e a recuperação de senha; quanto aos instrumentos, o Administrador pode consultá-los e excluí-los, mas não criá-los nem editá-los. O Administrador não é uma especialização do Psicopedagogo e não herda seus casos de uso: não gerencia aprendentes, não aplica instrumentos nem elabora relatórios, pois esses casos envolvem dados clínicos protegidos pelo sigilo profissional (Seção~\ref{sec:seguranca}). Cabem a ele o gerenciamento de usuários, a exclusão de instrumentos, a transferência de aprendentes entre profissionais e a consulta ao registro de auditoria sem os valores dos dados clínicos. O psicopedagogo, por sua vez, consulta os acessos realizados aos dados dos aprendentes sob sua responsabilidade, com os valores completos. A criação de uma nova versão é incluída no gerenciamento de instrumentos sempre que as questões de um instrumento já aplicado são alteradas, enquanto a finalização da aplicação e a impressão do relatório estendem, respectivamente, a aplicação do instrumento e a elaboração do relatório.

\begin{figure}[H]
    \centering
    \caption{Diagrama de casos de uso do sistema}
    \label{fig:casos_de_uso}
    \includegraphics[width=\textwidth,height=0.75\textheight,keepaspectratio]{casos_de_uso.png}
    \fonte{Elaborado pela autora (2026).}
\end{figure}

\subsection{Requisitos não funcionais}

Os requisitos não funcionais, apresentados no Quadro~\ref{qua:requisitos_nao_funcionais}, decorrem do fato de o sistema armazenar informações de crianças e priorizam segurança, privacidade, rastreabilidade e usabilidade. Seguindo a recomendação de expressá-los de forma mensurável (Seção~\ref{sec:fund_requisitos}), três deles receberam critérios de verificação: a usabilidade, avaliada por questionário com profissionais (RNF06); a acessibilidade, tomada a partir das diretrizes WCAG 2.2 \cite{w3c_wcag22_pt} e verificada pela ausência de violações detectáveis pela ferramenta automatizada axe-core \cite{deque_axe_core}, já que a conformidade plena depende de avaliação manual (RNF11); e a manutenibilidade, associada a uma cobertura mínima de testes automatizados no backend (RNF10).

\begin{quadro}[H]
\centering
\small
\caption{Requisitos não funcionais}
\begin{tabular}{p{2cm} p{13cm}}
\hline
\textbf{Código} & \textbf{Requisito não funcional} \\
\hline
RNF01 & O sistema deve utilizar autenticação para controlar o acesso aos usuários. \\
RNF02 & O sistema deve utilizar mecanismos de autorização para restringir o acesso às funcionalidades e informações, inclusive limitando a visualização de aprendentes e avaliações ao usuário responsável por cada registro. \\
RNF03 & Os dados dos aprendentes devem ser armazenados de forma segura, considerando os princípios de proteção de dados da LGPD. \\
RNF04 & O sistema deve manter um registro de auditoria que permita reconstituir quem criou, alterou, excluiu ou consultou cada dado, quando e com quais valores anteriores, inclusive nas operações realizadas pelo administrador. \\
RNF05 & O sistema deve preservar o histórico das versões e alterações realizadas nos instrumentos avaliativos. \\
RNF06 & A interface deve apresentar organização e usabilidade adequadas às atividades dos profissionais, verificada por meio de avaliação com usuários: média igual ou superior a 4,0 em escala de 1 a 5 nos aspectos de facilidade de uso, organização das informações e contribuição geral da plataforma. \\
RNF07 & A aplicação deve apresentar arquitetura em três camadas separadas (apresentação, aplicação e persistência), com o backend organizado em módulos independentes por domínio funcional (autenticação, usuários, aprendentes, avaliações, aplicações, relatórios, auditoria e dashboard), de modo que alterações em um módulo não exijam modificações nos demais. \\
RNF08 & Os dados devem possuir mecanismos de integridade e consistência durante seu armazenamento e manipulação. \\
RNF09 & O sistema deve possuir tratamento adequado de erros, evitando a exposição de informações internas da aplicação. \\
RNF10 & O sistema deve possuir cobertura de testes automatizados de no mínimo 70\% das linhas de código do backend, de modo a permitir a detecção de regressões durante a evolução do sistema. \\
RNF11 & A interface não deve apresentar violações dos critérios de sucesso dos níveis A e AA das WCAG 2.2 detectáveis pela ferramenta axe-core nas telas principais do sistema. \\
\hline
\end{tabular}
\fonte{Elaborado pela autora (2026).}
\label{qua:requisitos_nao_funcionais}
\end{quadro}